% Generated draft: requires mathematical and source review.
\documentclass[11pt]{article}
\usepackage[margin=25mm]{geometry}
\usepackage{fontspec}
\setmainfont{Latin Modern Roman}
\usepackage{amsmath,amssymb,mathtools,unicode-math}
\setmathfont{Latin Modern Math}
\usepackage{enumitem,needspace,xurl,hyperref}
\setlength{\parindent}{0pt}
\setlength{\parskip}{5pt}
\setlength{\emergencystretch}{3em}
\newcommand{\N}{\mathbb{N}}
\newcommand{\Z}{\mathbb{Z}}
\newcommand{\Q}{\mathbb{Q}}
\newcommand{\R}{\mathbb{R}}
\newcommand{\C}{\mathbb{C}}
\newcommand{\F}{\mathbb{F}}
\newcommand{\A}{\mathbb{A}}
\newcommand{\PP}{\mathbb{P}}
\newcommand{\E}{\mathbb{E}}
\newcommand{\eps}{\varepsilon}
\newcommand{\dd}{\,\mathrm{d}}
\begin{document}
\section*{UnsolvedMath ID: 30004543}
\textbf{Source problem code:} catalogue entry
\subsection{Definitions and mathematical statement}
Let
\[
S^6=\left\{x\in\R^7:\sum_{i=1}^{7}x_i^2=1\right\}
\]
carry its standard smooth structure. A complex-manifold structure compatible with this smooth structure is an atlas of charts to open subsets of \(\C^3\), with holomorphic transition maps, whose underlying smooth structure is the given one. Equivalently, it is an integrable almost-complex structure: a smooth bundle endomorphism \(J:TS^6\to TS^6\) satisfying \(J^2=-\operatorname{id}_{TS^6}\) and having vanishing Nijenhuis tensor
\[
N_J(X,Y)=[JX,JY]-J[JX,Y]-J[X,JY]-[X,Y]=0
\]
for all smooth vector fields \(X,Y\) on \(S^6\).

The principal question is:
\[
\text{Does the standard smooth manifold }S^6\text{ admit such an integrable }J?
\]
An almost-complex structure without the integrability condition does not answer this question.

The source also asks: which rational homotopy types occur as the underlying topological spaces of complex manifolds? In the simply connected setting, the rational homotopy type of a space \(X\) of CW type means the homotopy type of its rationalization \(X_{\Q}\); the rationalization map induces
\[
\pi_n(X)\otimes_{\Z}\Q\cong\pi_n(X_{\Q}),\qquad n\geq 2.
\]
This refers to the full rational homotopy type, not merely the list of rational homotopy groups. The source does not specify compactness, connectedness, simple connectedness, nilpotence, finite-type assumptions, or a fixed complex dimension for its broader question. In particular, it does not specify a convention for rational homotopy types outside the usual simply connected or nilpotent settings. These omissions leave the scope of that classification question unresolved; no particular restricted classification problem is selected here.
\subsection{Short English statement}
Does the standard smooth six-sphere admit complex coordinates with holomorphic transition maps, making it a complex manifold of complex dimension three? The source also asks which rational homotopy types complex manifolds can realize, but leaves the class of manifolds and the rational-homotopy convention insufficiently specified. The supplied record labels the problem open; this draft does not independently verify that status.
\subsection{Sources}
\begin{itemize}
\item Supplied local UnsolvedMath registry record, permanent ID \(30004543\), titled ``Complex Structures on the Six-Sphere and Rational Homotopy Types''; exact source problem code \texttt{catalogue entry}, in the Oberwolfach Reports Open Problems collection. No individual proposer is identified. The registry gives proposed year \(2020\).
\item Supplied report citation: \emph{Mini-Workshop: Almost Complex Geometry (2021)}, Oberwolfach Reports, pp.~1657--1691, DOI \texttt{10.4171/owr/2020/33}. \url{https://doi.org/10.4171/owr/2020/33}. The supplied background describes the report as documenting mathematical workshops and their open-problem discussions.
\item Supplied literature reference, labeled ``Chern study'': \url{https://arxiv.org/abs/1405.3405}.
\item Supplied literature reference, labeled ``Complex Manifolds 2025'': \url{https://doi.org/10.1515/coma-2025-0011}.
\item Literature assessment dated exactly \texttt{2026-08-22}: the supplied curation labels the problem open and states that existence of an integrable complex structure on the smooth six-sphere remains unresolved; known almost-complex structures and obstructions for restricted classes do not decide the general problem. This is supplied dated curation derived from the OWR record, not an independently verified literature assessment. Its time-sensitive status and cited claims require independent verification.
\item The registry labels the statement \texttt{corrected\_verified} and reports that notation, typography, omitted definitions, or duplicate clauses were repaired while mathematical content was preserved. This is a supplied editorial qualification, not mathematical verification performed for this draft.
\end{itemize}
\end{document}
